\documentclass[openright]{masterthesis}
\usepackage{FiraMono}

% enriicola's and pi's addings
\usepackage{array}
\usepackage{longtable}

\usepackage{graphicx}
\graphicspath{{assets/branding/}{assets/img/}}
\usepackage[font=scriptsize]{caption}
\usepackage{mwe}


\usepackage{soul}
\usepackage[dvipsnames]{xcolor}

\usepackage{graphicx}
\usepackage{hyperref}
\usepackage{tikz}
\usepackage{listings}
\usepackage{multicol}
\definecolor{chatIn}{HTML}{e2ffc9}
\definecolor{chatBack}{HTML}{efe6dd}

\definecolor{pastelBlue}{HTML}{AEAAE6}
\definecolor{pastelGreen}{HTML}{ACE8C0}
\definecolor{pastelYellow}{HTML}{FFF1AD}
\definecolor{pastelOrange}{HTML}{FFBA91}
\definecolor{pastelRed}{HTML}{E88E8E}
\definecolor{pastelGray}{HTML}{EDE4E0}

\newcommand{\comment}[1]{\textcolor{red}{#1}}

\lstset{basicstyle=\ttfamily\footnotesize, tabsize=3, breaklines=true, showspaces=false, showstringspaces=false, captionpos=b}

% macros
\newcommand*\javascript{JavaScript}
\newcommand*\nodejs{Node.js}

\begin{document}
\title{High Performance Ecoinformatics}
\author{Enrico Pezzano}
\advisor{Prof. Daniele D'Agostino\\Prof. Gabriele Casazza (DISTAV)}
\examiner{Prof. Matteo Dell'Amico}
\maketitle
~\newpage
\dedication{\textit{DEDICATION}}
~\newpage
\textbf{Abstract} Ecoinformatics is a modern method of doing science that enables scientists to generate new knowledge through innovative tools and approaches. Ecology is increasingly becoming a data-intensive science, involving the acquisition, preservation, and analysis of vast amounts of data from a variety of sources, including satellite and airborne remote sensing, instruments, sensors, and human observations. Ecologists are increasingly addressing larger-scale questions of both scientific and societal relevance, such as climate change and its consequences for ecosystems and related ecosystem services in the future. The topic of this thesis project stems from a research study conducted by the Department of Earth, Environmental, and Life Sciences (DISTAV, University of Genoa) that analyzes the shift in the distribution of natural habitats in the Alps due to future climate change to support decision-making processes regarding conservation planning and prioritization. The analysis is conducted using Species Distribution Models (SDMs), a dataset of approximately 3 million data points, and future climate change scenarios based on General Circulation Models (climate models). The statistical software R (https://www.r-project.org/) is used to generate future predictions of habitat distribution changes and biodiversity loss. The thesis candidate will work on improving the script and algorithmic functions, written in R, to streamline model analysis based on their knowledge and computer science background. This thesis project leverages interdisciplinary collaboration between DIBRIS and DISTAV to achieve better scientific results.   Depending on the candidate's interest and quality of results, the thesis can be upgraded with asterisk. 


~\newpage~\newpage

\tableofcontents

% \chapter{Temporary experimental notes}

\section{Experimental question}

The Leonardo experiments investigate how BIOMOD2 projection storage affects runtime and memory use. Each species is calibrated with five algorithms (GLM, GBM, ANN, FDA and MAXNET) and then projected over the current raster and eight future climate scenarios. Each raster contains approximately 64 million cells, so retaining many model projections simultaneously can exhaust the approximately 494~GB available on a DCGP node.

The immediate objective was to complete one full production-scale run for \emph{Achillea atrata}. This provides a runtime and resource baseline before considering a refactor or launching the full 167-species campaign.

\section{Run inventory}

Table~\ref{tab:run-inventory} lists every run or campaign for which the current project notes contain an identifier. Entries marked ``not recorded'' are intentionally left incomplete. Rows use different species, pseudo-absence settings, model counts, code revisions and storage modes, so the table is an inventory rather than a controlled comparison. In the storage column, T/T, F/T, T/F and F/F abbreviate \texttt{keep.in.memory}/\texttt{do.stack}; ``default T/T'' means that both arguments were omitted and BIOMOD2 defaults applied.

{\footnotesize
\setlength{\tabcolsep}{2pt}
\begin{longtable}{>{\raggedright\arraybackslash}p{2.1cm}>{\raggedright\arraybackslash}p{2.0cm}>{\raggedright\arraybackslash}p{1.5cm}>{\raggedright\arraybackslash}p{1.7cm}>{\raggedright\arraybackslash}p{2.0cm}>{\raggedright\arraybackslash}p{4.2cm}}
\caption{Chronological inventory of documented Leonardo runs and campaigns.}\label{tab:run-inventory}\\
\hline
Run/ref. & Scope & CPU & Mode & State/time & Recorded observation \\
\hline
\endfirsthead
\hline
Run/ref. & Scope & CPU & Mode & State/time & Recorded observation \\
\hline
\endhead
\texttt{52004b2}, \texttt{354e415} & legacy tests & $\geq16$ & not recorded & failed / up to 12 h & OOM during ensemble forecasting; some runs also reached the 12-hour limit. \\
\texttt{bc71039} & legacy single species & 1 & not recorded & failed / 12 h & Time limit reached during future scenario 2 of 8. \\
Unidentified early run & 9-row \emph{A. atrata} test & 1 & not recorded & completed / 7:01 & Toy parameters; 24 surviving models. \\
\texttt{e622491} & small test input & 1 & default T/T & completed / 50.3 h & Current projection and eight future scenarios completed. \\
\texttt{b26fbd3} & 1,000-row Agrostis test & 1 & default T/T & completed / 17.2 h & Toy parameters; approximately 30 models and eight future scenarios. \\
Preliminary split versus monolithic series & small input, eight runs & 1 & not applicable & completed / 25:51--26:28 & Four split and four monolithic runs; MaxRSS range 40.3--48.3 GiB. \\
Controlled split versus monolithic series & small input, ten runs & 1 & not applicable & completed / 25:59--27:12 & Five split and five monolithic runs; MaxRSS range 47.04--48.16 GiB. \\
\texttt{47333938} & \emph{A. capillaris}, 170,701 rows & 4 & default T/T & completed / 4:23 & Ten pseudo-absences; 24 surviving models; peak approximately 326 GB. \\
\texttt{47467973} & \emph{A. capillaris}, 170,701 rows & 4 & default T/T & completed / 11:50 & 10,000 pseudo-absences; 30 valid models; peak approximately 359 GB. \\
\texttt{47510573\_1} & \emph{P. erecta}, 167,345 rows & 4 & default T/T & completed / 11:33 & Peak approximately 360 GB; timing file present. \\
\texttt{47510573\_2} & \emph{G. anisophyllon}, 5,936 rows & 4 & default T/T & cancelled / $\sim$18 h & Projection entered a SIGPIPE retry loop. \\
\texttt{47510573\_3} & \emph{F. glauca}, 332 rows & 4 & default T/T & cancelled / $\sim$18 h & Projection entered a SIGPIPE retry loop. \\
\texttt{47574797\_2} & \emph{G. anisophyllon}, 5,936 rows & 1 & default T/T & failed / 3:12 & Ensemble output reported \texttt{writeRaster path does not exist}. \\
\texttt{47574797\_3} & \emph{F. glauca}, 332 rows & 1 & default T/T & completed / $\sim$23 h & Sequential small-species rerun. \\
\texttt{47593184\_*} & array, scope not recorded & not recorded & not recorded & mixed & Slurm summary recorded a mixed state and maximum signal 9; task detail was not preserved in these notes. \\
\texttt{48075655} & 167-species array & 1 or 4 & default T/T & first five OOM / 13--17 h & First five tasks failed during future projection; remaining tasks were cancelled. \\
\texttt{48236130\_1} & \emph{A. atrata} & 1 & F/F & interrupted & Stopped before projection when the test species was changed. \\
\texttt{48238919\_3} & \emph{A. capillaris} & 4 & F/F & interrupted / $\sim$22 h & Reached future scenario 2 of 8; no final completion result. \\
\texttt{48325677\_3} & \emph{A. capillaris} & 4 & not recorded & failed / 8.45 h & Parallel ensemble projection reported that an output file already existed. \\
\texttt{48325677\_4} & \emph{A. rupestris} & 1 & not recorded & not recorded & Last preserved status was running during pseudo-absence selection. \\
\texttt{48325677\_5} & \emph{A. fissa} & 1 & not recorded & not recorded & Last preserved status was running during pseudo-absence selection. \\
Unidentified 2026-07-05 array & first five SJF species & not recorded & not recorded & failed & With a 100-GB request, all five tasks were recorded as OOM during current ensemble forecasting. \\
\texttt{49507005--}\newline\texttt{49507067} & 56-wave, 167-species campaign & 5 & F/F & cancelled & No valid completion was found in the later audit; first-wave live MaxRSS was approximately 261--316 GiB. \\
\texttt{49842976\_1} & \emph{A. atrata} & 2 & F/T & failed / 05:44:00 & Fork allocation failed during the first future clamping mask; MaxRSS 413.22 GiB. \\
\texttt{49843592\_1} & \emph{A. atrata} & 4 & F/T & failed / 03:57:25 & Same observed failure point; MaxRSS 413.80 GiB. \\
\texttt{49844162\_1} & \emph{A. atrata} & 4 & F/F & completed / 22:33:50 & MaxRSS 260.94 GiB; eight future scenarios and \texttt{\_SUCCESS}. \\
\texttt{51485472\_1} & \emph{A. atrata} & 4 & T/F & completed / 22:21:40 & MaxRSS 264.69 GiB; 1,368 non-empty output files. \\
\texttt{51494635\_1} & \emph{A. atrata} & 6 & F/F & completed / 17:50:56 & MaxRSS 258.95 GiB; 1,368 non-empty output files. \\
\texttt{51485581\_1} & \emph{A. atrata} & 8 & F/F & completed / 15:05:35 & MaxRSS 278.67 GiB; 1,368 non-empty output files. \\
\texttt{51738981\_1} & \emph{A. atrata} & 16 & F/F & completed / 10:21:30 & Final MaxRSS and output validation not recorded. \\
\texttt{51739004\_1} & \emph{A. atrata} & 16 & T/F & OOM / 02:57:18 & One \texttt{oom\_kill} during first future scenario; MaxRSS 472.57 GiB. \\
\texttt{51739048\_1} & \emph{A. atrata} & 32 & F/F & OOM / 01:39:58 & Twelve \texttt{oom\_kill} events during current projection; MaxRSS 481.16 GiB. \\
\texttt{51756264\_1} & \emph{A. atrata} & 12 & F/F & time limit / 09:30:20 & Incomplete under the imposed 9:30 limit. \\
\texttt{51756286\_1} & \emph{A. atrata} & 12 & T/F & time limit / 09:30:18 & Incomplete under the imposed 9:30 limit. \\
\hline
\end{longtable}
}

\section{Meaning of the storage options}

BIOMOD2 defaults both \texttt{keep.in.memory} and \texttt{do.stack} to \texttt{TRUE}. When \texttt{do.stack} is enabled, the individual model projections are combined into a single multilayer raster. When \texttt{keep.in.memory} is enabled, that projection object is also retained in the R object returned by BIOMOD2.

The disk-backed configuration sets both options to \texttt{FALSE}. It does not eliminate RAM use: R still needs memory for environmental rasters, fitted models, worker processes, clamping masks and temporary buffers. The difference is that each model projection is written to its own file and the returned BIOMOD2 object retains links to those files instead of the complete raster stack. This prevents all projection layers from accumulating in memory at the same time, at the cost of additional filesystem I/O.

The tested hybrid configuration used \texttt{keep.in.memory = FALSE} but \texttt{do.stack = TRUE}. It avoided retaining the final stack in the returned object, but BIOMOD2 still had to construct the multilayer stack during projection. Jobs \texttt{49842976\_1} and \texttt{49843592\_1} both stopped during first-scenario clamping-mask construction after reporting a fork allocation failure.

\section{Evidence from earlier successful runs}

The Git history and preserved logs confirm that the earlier successful scripts did not pass either storage option to \texttt{BIOMOD\_Projection}. They therefore used the BIOMOD2 defaults, \texttt{TRUE/TRUE}. Those successes are real, but they used substantially fewer model combinations than the current production configuration.

Several examples document the difference. The early sequential run at commit \texttt{e622491} used two pseudo-absence repetitions and two cross-validation repetitions on the small test input. Later Agrostis experiments used three pseudo-absence repetitions, two cross-validation repetitions and five algorithms, producing approximately 30 models.

Job \texttt{47467973} completed with 10,000 pseudo-absences, 30 valid models, four workers, a wall time of 11:50 and a peak near 359~GB. Job \texttt{47510573\_1} similarly completed \emph{Potentilla erecta} in 11:33 with a peak near 360~GB. Job \texttt{47574797\_3} completed \emph{Festuca glauca} sequentially in approximately 23 hours.

At least one successful historical run used 10,000 pseudo-absences. The documented configurations generated different numbers of models: roughly 30--35 in the historical runs, 125 in the present baseline (five pseudo-absence repetitions by five cross-validation repetitions by five algorithms), and up to 250 in the later ten-repetition configuration. Because pseudo-absence repetitions, cross-validation repetitions, valid-model counts and storage settings differ simultaneously, these historical runs are not a controlled storage comparison.

\section{First full-array and disk-backed tests}

The first direct attempt to process all 167 species was array \texttt{48075655\_[1-167]}, submitted at the end of June 2026 with a three-day wall-time request. Each array index selected one species, while the long-production QOS allowed at most three one-node tasks to run concurrently. The array initially remained pending because its requested duration overlapped a DCGP maintenance reservation from 2026-06-30 at 09:00 until 2026-07-01 at 09:00.

After maintenance, the first five tasks started in groups of three. All five were killed for out-of-memory conditions after approximately 13--17 hours. The failures covered \emph{Achillea atrata} with 1,475 occurrence rows, \emph{Achillea clusiana} with 174, \emph{Agrostis capillaris} with 170,701, \emph{Agrostis rupestris} with 4,464 and \emph{Alchemilla fissa} with 1,428. Four tasks used one projection worker under the contemporary small-species workaround, whereas \emph{Agrostis capillaris} used four. The observed failures therefore included both small and large occurrence counts and both one- and four-worker tasks.

The historical configuration used ten pseudo-absence repetitions, five cross-validation repetitions and five algorithms, allowing approximately 250 model runs per species. Because the script omitted the projection-storage arguments, BIOMOD2 used \texttt{TRUE/TRUE}. The failures occurred during future projection while the script was using the default in-memory stacked representation. With explicit approval, the remaining array was cancelled to avoid repeating the same failure across 167 tasks. Commit \texttt{434b6a8} then set both storage options to \texttt{FALSE} for current and future projections.

Job \texttt{48236130\_1} began validation of that change on \emph{Achillea atrata} with one worker. It was intentionally stopped before reaching projection when the test target was changed to the larger, four-worker \emph{Agrostis capillaris}. The replacement, job \texttt{48238919\_3}, reached future scenario 1 of 8 after approximately 14 hours 50 minutes and scenario 2 after approximately 20 hours 15 minutes. At a live check it had run for about 22 hours and was still writing projection rasters without an OOM. It was nevertheless interrupted before completion when a different three-species experiment was launched, so it provides evidence of progress but not a completed validation.

Commit \texttt{6c36715} introduced an experimental test worker and configurable switches in \nolinkurl{scripts/sbatch.sh}. The experimental script reused already loaded calibration rasters, released the complete occurrence table after selecting one species, configured Terra's memory fraction and temporary directory, and added output-reuse controls. Array \texttt{48325677}, with tasks 3--5, then ran three species on three DCGP nodes. The jobs initially appeared inactive while BIOMOD2 selected random pseudo-absences, but live CPU accounting showed approximately one core fully occupied. Subsequent profiling identified repeated scans of the 64-million-cell raster during ten pseudo-absence repetitions as a multi-hour formatting bottleneck.

Task \texttt{48325677\_3} later failed after 507 minutes (8.45 hours) during parallel \nolinkurl{BIOMOD_EnsembleForecasting}. Multiple workers attempted to write the same disk-backed projection file, and Terra reported that the file already existed. This observed error differed from the earlier OOM: multiple workers reported a collision on the same disk-backed projection path. Later commits continued from this experimental branch with deterministic seeds, overwrite handling, shortest-job-first ordering and success markers; the current production baseline instead resides in \nolinkurl{R/base/base_sequential_analysis.R}.

\section{Controlled storage comparison}

Three isolated jobs processed task 1, \emph{Achillea atrata}, while reserving the full node memory. Job \texttt{49842976\_1} tested hybrid storage with two workers, job \texttt{49843592\_1} tested the same mode with four workers, and job \texttt{49844162\_1} tested disk-backed storage with four workers. CPU allocation and billing matched the requested worker count because \texttt{--exclusive} was not used.

Both hybrid jobs failed at the same point: construction of the clamping mask for the first future scenario. The R runtime reported that \texttt{mcfork()} could not allocate memory. The two-worker run failed after 05:44:00 with a Slurm MaxRSS of 413.22~GiB; the four-worker run reached the same failure after 03:57:25 with 413.80~GiB. Slurm did not label either job \texttt{OUT\_OF\_MEMORY}; the application log reported \texttt{mcfork(): unable to fork, possible reason: Cannot allocate memory}.

The four-worker disk-backed job completed successfully. Its wall time was 22:33:50, and Slurm reported a MaxRSS of 260.94~GiB. The \texttt{\_SUCCESS} marker records completion at 2026-07-20 18:37:21. Both timing files are present and list all eight future scenarios. The species output contains 1,364 non-empty files occupying approximately 7.2~GB, including 125 individual models, two ensembles and all current and future projection directories.

Most of the successful run was spent in projection. Formatting required 4,527.93~s (75.47 minutes; 1.26 hours), individual modeling 248.85~s (4.15 minutes; 0.07 hours), ensemble modeling 137.30~s (2.29 minutes; 0.04 hours), current projection 7,047.56~s (117.46 minutes; 1.96 hours), current ensemble projection 1,421.92~s (23.70 minutes; 0.40 hours), and future projections 67,818.98~s (1,130.32 minutes; 18.84 hours). No fatal application error appears in the log. Non-fatal GLM, MAXNET and ensemble-metric warnings are present and remain to be audited.

\section{Measured worker comparison}

The following runs compare observed wall time and memory for one species. ``Disk-backed'' should not be interpreted as ``no RAM'': job \texttt{49844162\_1} used 260.94~GiB. In this representation, completed projection layers are persisted separately and referenced from disk instead of being accumulated into one multilayer stack.

Three further single-species jobs from commit \texttt{88e03c0} completed on 2026-08-02. Each reserved 494000~MB, processed task 1 and used an isolated output directory.

Job \texttt{51485472\_1} ran four workers with \texttt{TRUE/FALSE} on \texttt{lrdn4939}. Job \texttt{51494635\_1} ran six workers with \texttt{FALSE/FALSE} on \texttt{lrdn4454}. Job \texttt{51485581\_1} ran eight workers with \texttt{FALSE/FALSE} on \texttt{lrdn4946}. Ensemble stages remained capped at two workers.

The four- and eight-worker jobs were initially submitted with the script's four-day limit. They remained pending even though 127 DCGP nodes appeared idle because a maintenance reservation begins on 2026-08-03 at 08:00 and ends on 2026-08-17 at 18:00. After explicit approval, their limits were reduced in place to 30 hours and both started at approximately 16:50. The six-worker job was submitted directly with the same limit and started at 16:54:58. All three completed before maintenance with Slurm state \texttt{COMPLETED} and exit code zero.

The four-worker \texttt{TRUE/FALSE} job took 80,500~s (1,341.67 minutes; 22.36 hours) and reached 264.69~GiB MaxRSS. The six-worker disk-backed job took 64,256~s (1,070.93 minutes; 17.85 hours) and reached 258.95~GiB. The eight-worker disk-backed job took 54,335~s (905.58 minutes; 15.09 hours) and reached 278.67~GiB. These values use the Slurm batch-step MaxRSS, which captures the whole job cgroup and is therefore preferred over the lower per-process maximum reported by GNU \texttt{time}.

Relative to the original four-worker \texttt{FALSE/FALSE} baseline of 81,230~s (1,353.83 minutes; 22.56 hours), enabling \texttt{keep.in.memory} while retaining separate files saved only 730~s (12.17 minutes; 0.20 hours), or 0.9\%, while increasing peak memory by 3.75~GiB. Only one run was performed for each configuration, so no run-to-run variance was measured. The six-worker run achieved a speedup of 1.264 and reduced wall time by 20.9\%. The eight-worker run achieved a speedup of 1.495 and reduced wall time by 33.1\%. Moving from six to eight workers produced a further speedup of 1.183 and a 15.4\% wall-time reduction.

The timing files report the following phase durations.

With four workers and \texttt{TRUE/FALSE}, formatting took 4,490.02~s (74.83 minutes; 1.25 hours), current projection 6,988.68~s (116.48 minutes; 1.94 hours), current ensemble projection 1,467.62~s (24.46 minutes; 0.41 hours), and the complete future loop 67,139.91~s (1,119.00 minutes; 18.65 hours).

At six workers these phases took 4,600.66~s (76.68 minutes; 1.28 hours), 5,132.52~s (85.54 minutes; 1.43 hours), 1,451.78~s (24.20 minutes; 0.40 hours), and 52,735.75~s (878.93 minutes; 14.65 hours).

At eight workers they took 4,531.17~s (75.52 minutes; 1.26 hours), 4,057.93~s (67.63 minutes; 1.13 hours), 1,448.78~s (24.15 minutes; 0.40 hours), and 43,996.87~s (733.28 minutes; 12.22 hours).

Across these runs, formatting remained between 4,490.02 and 4,600.66~s (74.83--76.68 minutes; 1.25--1.28 hours), while current ensemble projection remained between 1,448.78 and 1,467.62~s (24.15--24.46 minutes; 0.40--0.41 hours).

Slurm recorded 70.68 CPU-hours for the four-worker run, 72.95 for six workers and 72.51 for eight. These correspond to average utilization of 3.16, 4.09 and 4.80 cores, or approximately 79.0\%, 68.1\% and 60.0\% of the allocated CPUs. Allocated CPU-hours were 89.44, 107.09 and 120.74 respectively.

All three output directories contain 1,368 files, no empty files, both timing tables and a non-empty \texttt{\_SUCCESS} marker. Their sizes are 7.3--7.4~GB. After normalizing BIOMOD2's run-specific numeric identifier, their relative file structures are identical. No fatal error, OOM or fork failure appears in the logs. The repeated GLM probability, integer-overflow, missing-metric and optional \texttt{cito} warnings remain scientific-review items rather than execution failures.

To check for accidental output reuse, the worker and outputs were inspected after completion. The worker deletes the complete species directory before formatting. None of the 1,368 files in any output predates its job start, and no log contains a skip or pre-existing-file message. Each log records all eight future scenarios and 1,143 projection writes, corresponding to 125 individual models plus two ensembles across the current environment and eight future environments. The approximately 71--73 CPU-hours recorded per run provide independent evidence of full recomputation.

These measurements concern one species and one run per configuration. No configuration for the 167-species campaign is selected in these notes.

\subsection{High-worker follow-up}

Three further isolated tests were submitted from commit \texttt{88e03c0} on 2026-08-02. Because DCGP maintenance begins at 08:00 on 2026-08-03, each requests a 14-hour limit and must finish shortly before the reservation. The 14-hour limit was selected so that the requested interval ended before maintenance; completion within that limit was not known at submission. The 32-worker run records the behavior of a larger worker count under the same node-memory request.

Job \texttt{51738981\_1} ran 16 workers with \texttt{FALSE/FALSE} on \texttt{lrdn4795} and wrote to \nolinkurl{data/output_disk_sp1_16cpu}. A Slurm notification recorded completion with exit code zero after 37,290~s (621.50 minutes; 10.36 hours). At an earlier live check after 15,943~s (265.72 minutes; 4.43 hours), it was processing future scenario 3 of 8 and had reached 464.77~GiB MaxRSS. Leonardo entered maintenance before final accounting and output validation could be collected. Final MaxRSS and the output contents of job \texttt{51738981\_1} will therefore be collected and validated after maintenance.

Job \texttt{51739004\_1} ran 16 workers with \texttt{TRUE/FALSE} on \texttt{lrdn3936} and wrote to \nolinkurl{data/output_keepmem_sp1_16cpu}. Slurm recorded \texttt{OUT\_OF\_MEMORY}, exit code \texttt{0:125}, and an elapsed time of 10,638~s (177.30 minutes; 2.96 hours). Batch-step MaxRSS was 472.57~GiB. During individual-model projection for the first future scenario, one scheduled worker did not return; Slurm reported one \texttt{oom\_kill}. The subsequent ensemble call stopped because the projection result did not contain every requested model.

Job \texttt{51739048\_1} ran 32 workers with \texttt{FALSE/FALSE} on \texttt{lrdn3756} and wrote to \nolinkurl{data/output_disk_sp1_32cpu}. Slurm recorded \texttt{OUT\_OF\_MEMORY}, exit code \texttt{0:125}, and an elapsed time of 5,998~s (99.97 minutes; 1.67 hours). Batch-step MaxRSS was 481.16~GiB. During current individual-model projection, 12 scheduled workers did not return and Slurm reported 12 \texttt{oom\_kill} events. The subsequent ensemble call stopped because models were missing from the projection result.

All three output directories were absent immediately before submission. Initial logs confirmed matching Slurm CPU allocations, BIOMOD worker counts and storage flags; ensemble stages remained capped at two workers. At the first check all three jobs were running and together occupied the project's three-node allowance. Two later reached Slurm state \texttt{OUT\_OF\_MEMORY}; the 16-worker \texttt{FALSE/FALSE} job later completed according to its Slurm notification.

After those two failures released their nodes, two 12-worker jobs were submitted with new isolated directories. Job \texttt{51756264\_1} uses \texttt{FALSE/FALSE} on \texttt{lrdn4863} and writes to \nolinkurl{data/output_disk_sp1_12cpu}. Job \texttt{51756286\_1} uses \texttt{TRUE/FALSE} on \texttt{lrdn4866} and writes to \nolinkurl{data/output_keepmem_sp1_12cpu}. Their logs confirm 12 BIOMOD workers, two ensemble workers and matching Slurm allocations.

Each 12-worker job had a limit of 34,200~s (570 minutes; 9.50 hours). They started at 22:19:55 and 22:20:27 and reached \texttt{TIME\_LIMIT} after 34,220~s (570.33 minutes; 9.51 hours) and 34,218~s (570.30 minutes; 9.51 hours), respectively. These elapsed times describe the imposed limits rather than the complete workload runtime.

\subsection{Observed memory mechanism and campaign storage}

The container uses BIOMOD2 4.3.4.5 and Terra 1.9.11. Terra reported a default memory fraction of 0.5, no absolute memory maximum and a temporary directory under \texttt{/tmp}. This limit applies within each process; it does not impose one aggregate allowance across all BIOMOD2 workers.

Code inspection identified a distinct projection path for MAXNET. The other four algorithms in this experiment delegate raster prediction to Terra, which can process raster blocks. The MAXNET method first converts the complete environmental raster with \texttt{as.points()} and then \texttt{as.data.frame()} before prediction. The experiment raster contains 63,951,097 cells and five variables. Each species has 25 MAXNET models among 125 individual models, so concurrent MAXNET workers can hold multiple complete-raster representations at once.

A local diagnostic used one stored MAXNET model and a 230,509-cell raster crop containing 226,775 complete cells. The existing BIOMOD2 path was compared with a Terra prediction path capable of block processing, using the same fitted model. The two outputs had identical missing-value positions and a maximum absolute difference of zero on the compared cells. This diagnostic does not yet validate complete-raster runtime, peak memory or all 25 MAXNET models.

With \texttt{do.stack = FALSE}, the inspected BIOMOD2 code returns links to individual files. The \texttt{keep.in.memory} branch that stores projection values is inside the \texttt{do.stack} branch; with separate files, setting \texttt{keep.in.memory = FALSE} additionally frees the returned projection object. The temporary memory used within each concurrent model prediction remains separate from these output-retention settings.

The project filesystem reported approximately 1.0~TiB total, 225~GiB used and 800~GiB available at the check. One completed species occupied 7.4~GiB: 6.6~GiB in current and future individual-model projection directories, 700~MiB in ensemble projection directories, and 148~MiB in model and BIOMOD metadata directories. These are measurements for one \emph{A. atrata} output, not measurements of every species. If all 167 outputs had the same size, retaining every file would require approximately 1.21~TiB. Retaining only the measured 700~MiB of ensemble directories and 148~MiB of model and BIOMOD metadata directories would require approximately 138~GiB before small top-level files are added. This is a disk-space estimate, not a RAM measurement. Clamping masks reside inside the individual-projection directory structure, so preserving them while removing model projection rasters requires a selective retention rule and a new size measurement.

The \texttt{\_SUCCESS} marker states that the R script reached its final write step; it does not establish that every output has been independently validated or that linked individual projections are no longer required. Removing those rasters would prevent later inspection of individual-model maps and would invalidate BIOMOD objects that link to them. It would also require recomputation before changing ensemble definitions or thresholds. No automatic deletion policy has been applied.

\subsection{Measurement commands and interpretation}

Final Slurm state, exit code, elapsed time and memory were queried with \texttt{sacct} for the relevant job IDs and their batch steps. Live state and deadlines were queried with \texttt{squeue}, while \texttt{sstat} supplied live batch-step MaxRSS during running jobs. The reported values with a \texttt{K} suffix were converted to GiB by dividing by 1,048,576. In these one-task jobs, the Slurm batch-step accounting includes the R process and its forked descendants and is preferred to the smaller single-process maximum written by GNU \texttt{time}. Final Slurm accounting for job \texttt{51738981\_1} was unavailable when maintenance closed the login session.

Log tails located the operation immediately preceding each failure. Searches for worker-count and storage debug lines confirmed that command-line allocations matched the R settings. Searches for future-scenario markers identified progress through the eight scenarios. The 32-worker log reported 12 workers that did not return followed by 12 Slurm OOM-kill events; the 16-worker \texttt{TRUE/FALSE} log reported one missing worker and one OOM-kill event. The later BIOMOD ensemble errors listed missing model projections and occurred after those worker losses.

Container inspection used \texttt{singularity exec} to print BIOMOD2 and Terra versions and \texttt{terraOptions()}. Local source inspection then followed \texttt{BIOMOD\_Projection} into the algorithm-specific prediction methods. The MAXNET method converts the raster twice. It first uses \texttt{as.points()}. It then uses \texttt{as.data.frame()}. An attempted direct process listing on the compute nodes was rejected by Leonardo authentication, so no per-process RSS claim is based on that attempt.

Disk allocation was measured with \texttt{du}, whereas filesystem capacity was read with \texttt{df}. These commands measure disk blocks, not resident memory. Separate \texttt{du} measurements produced three totals for the completed eight-worker output. Individual projection directories occupied 6.6~GiB. Ensemble directories occupied 700~MiB, while model and BIOMOD metadata directories occupied 148~MiB. \texttt{gdalinfo} reported LZW compression on a sampled individual GeoTIFF.

\subsection{Partition and QOS notes}

All three submissions use \texttt{scripts/sbatch.sh}. Command-line options override only the CPU count, time limit and experiment-specific environment variables. The script continues to provide the account, partition, QOS, memory request, logging and mail configuration.

\texttt{dcgp\_usr\_prod} is the general-purpose CPU partition used by these experiments. Its normal partition time limit is one day. Different QOS choices modify scheduling limits for projects authorized to use them.

\texttt{dcgp\_qos\_lprod} is the long production QOS selected by the script. It permits wall times up to four days and, for this project, a maximum aggregate allocation of three nodes, 336 CPUs and 1482000~MB. During the comparison, the three one-node tests filled the observed three-node project allowance.

\texttt{dcgp\_qos\_bprod} is another production QOS available to the project. It has no separate wall-time extension in the reported configuration, so jobs inherit the one-day limit of \texttt{dcgp\_usr\_prod}.

\texttt{dcgp\_qos\_dbg} is the debug QOS. It permits jobs up to 30 minutes and reports a project limit of two nodes, 224 CPUs and 988000~MB.

\texttt{boost\_usr\_prod} is the production partition for GPU nodes equipped with four A100 devices. Its reported partition time limit is one day. \texttt{lrd\_all\_serial} provides two serial nodes with a four-hour limit, while \texttt{lrd\_all\_viz} provides visualization nodes with Quadro GPUs and a 12-hour limit. These alternatives describe different hardware and scheduling targets rather than interchangeable queues for the same DCGP workload.

The first three worker-count tests are complete. In the 16/16/32-worker follow-up, two jobs reached \texttt{OUT\_OF\_MEMORY} and the 16-worker \texttt{FALSE/FALSE} job completed according to its Slurm notification. Two additional 12-worker jobs reached their shorter time limits. The observed long-QOS allowance permits at most three one-node tasks concurrently. No full-campaign configuration or submission has been selected. RAM from separate nodes is not combined by these jobs; each array task processes one species independently. Slurm mail remains configured for array-task completion, failure and time-limit notifications.

\section{Full-campaign baseline objective}

The campaign-level objective is to measure the execution of the complete 167-species dataset with a fixed, documented version of the R workflow. This measurement is distinct from selecting the largest worker count that can fit on one node. Its record must include campaign makespan, the sum of per-species elapsed times, total CPU-hours, completed and failed species, per-phase timings, MaxRSS and output validation. These quantities distinguish elapsed calendar time under three-node concurrency from the aggregate time that the same species tasks would occupy sequentially.

The investigation is divided into four stages. First, live job accounting and code inspection identify the allocations that grow with each projection worker. Second, the unchanged R workflow and its outputs define the baseline independently of later implementation changes. Third, the campaign submission must allow individual species to complete or fail without discarding the measurements from other array indices. Fourth, an optimized implementation can be compared with the recorded baseline using the same models, raster cells, scenarios and output checks.

The memory evidence narrows the projection bottleneck to concurrent complete-raster representations rather than CPU allocation alone. As detailed above, the 100 GLM, GBM, ANN and FDA models use Terra prediction paths that can process blocks, whereas the 25 MAXNET models convert all 63,951,097 cells to points and then to a data frame. Multiple concurrent MAXNET workers therefore materialize coordinates, five predictor columns and prediction intermediates at the same time.

The local MAXNET comparison using the block-capable Terra path is retained as a diagnostic result rather than a production change. It used an existing fitted model and the same predictor rasters without consuming a Leonardo compute node. Equality on 226,775 valid cells shows that this implementation path can preserve the sampled prediction values, but complete-raster memory, runtime and equality across all models and scenarios remain unmeasured. No 167-species submission has been made on the basis of this prototype.

\subsection{Occurrence rows and raster blocks}

The occurrence CSV contains 2,583,359 rows across 167 species. \emph{A. atrata} contributes 1,475 rows, while the observed species counts range from 43 to 170,701. The number of occurrence rows affects formatting and model fitting, but every completed species still requests 125 model projections over the same current raster and eight future rasters. Therefore the 6.6~GiB measured for one species cannot be multiplied in proportion to occurrence rows. GeoTIFF compression also depends on predicted values and missing-value patterns. A size-versus-occurrence regression requires completed outputs from several species under the same model count, raster set and storage configuration; the available controlled comparison contains only \emph{A. atrata}.

Splitting occurrence rows within one species would change pseudo-absence selection, cross-validation and fitted models, so the partial results would not reconstruct the existing analysis. Grouping complete species into larger tasks would preserve the analysis but would not reduce projection memory and would weaken per-species failure isolation. Raster blocking is different: one fitted model predicts consecutive spatial row ranges or tiles, writes each block and combines them into the same output raster. It changes the memory schedule rather than the occurrence data or fitted model.

\subsection{Experimental MAXNET variant}

A separate temporary experimental entry point, \nolinkurl{R/tmp/base_sequential_analysis_blockwise.R}, loads \nolinkurl{R/tmp/maxnet_blockwise_override.R} before the unchanged baseline script. The override registers a more specific MAXNET \texttt{predict()} method that delegates numeric raster prediction to Terra. It explicitly rejects categorical rasters and scaled models, neither of which is used by the current workflow. The original \nolinkurl{R/base/base_sequential_analysis.R} is unchanged.

The runnable check \nolinkurl{tests/test_maxnet_blockwise.R} compares the original and experimental paths using a stored MAXNET model. It verifies identical missing-value positions, zero maximum absolute difference, identical integer-scaled values and unchanged data-frame prediction behavior on 226,775 valid raster cells. This check passed locally in the production container. It does not replace a full-raster Leonardo comparison of runtime, MaxRSS and output files.

\subsection{Deferred validation and maintenance access}

A read-only post-job validation plan was prepared for job \texttt{51738981\_1}. It would request one CPU, 4~GiB and 30 minutes, report final Slurm accounting, and check \texttt{\_SUCCESS}, timing and evaluation tables, total and empty files, 225 MAXNET rasters, eight future timing rows, eight logged scenarios, 1,143 projection writes, fatal log markers, disk use and files predating the job start. It was not submitted.

Leonardo broadcast that all user sessions would close at 09:30 on 2026-08-03 and then disconnected the active login. A later connection attempt through \nolinkurl{scripts/leogin.sh} reached the CINECA SSO flow but timed out before authentication, so it did not establish whether a login node was reachable after authentication. The observed Slurm reservation extended to 2026-08-17 18:00. The CINECA maintenance notice described an expected compute-node outage from August 3 to August 7 and stated that normal production was not guaranteed through August 14; access and filesystems could also be unavailable. Consequently, production jobs cannot run during the compute outage, but queue submission may become possible before August 17 if login, Slurm and filesystem services return. No validation or campaign job is currently queued.

\input{Chapters/1_Introduction}
% \input{Chapters/IntroductionAndMotivation} % Your work
\chapter{Background}

\comment{real world and eventually related work}

\section{Ecoinformatics}

\section{Biological study context}

An internal DISTAV methods document describes the biological study supported by the computational workflow. The study examines how climate change may alter the suitability of Alpine grassland habitats. It considers ten dry, alpine and subalpine habitat types from the EUNIS classification and 167 diagnostic plant species. Occurrence records are associated with environmental predictors at a spatial resolution of approximately 1~km. The predictors describe climate, soil and topographic heterogeneity. Species Distribution Models combine five modelling techniques and ensemble forecasts to estimate present and future habitat suitability.

The biological protocol also defines the analyses that follow model projection: changes in suitability, spatial shifts, habitat-level hotspot and cold-spot maps, and the relationship between suitability changes and species tolerance. This document provides the scientific context for the computing work, but it is not treated as the exact specification of every run. Dataset sizes, model repetitions and available climate scenarios changed during the project; the computational chapters therefore report the configuration recorded by the corresponding scripts and run manifests.

\section{Data used}

\section{Related work}

% \input{Chapters/Background}        % Work done by others!
% \input{Chapters/RelatedWork} % Work done by others!
\chapter{Code as is}
\comment{Existing workflow and computational problem.}

\section{R}

\comment{BEGIN TODO: Container, workflow and Leonardo}

The workflow received from DISTAV was a single R script written for a Windows workstation. It used absolute paths, loaded both \texttt{raster} and \texttt{terra}, and created a ten-worker \texttt{doParallel} cluster. The filename \nolinkurl{ensamble_modelling_no_parallel.R} is therefore misleading: model fitting and projection both used \texttt{nb.cpu = 10}. The script selected one species by a fixed index and limited that species to 100,000 occurrence rows. It also contained path-dependent parsing and a syntax error in the future ensemble call. The historical file is preserved without correction so that it remains a record of the received code.

\section{BIOMOD2}

BIOMOD2 performs the scientific stages of the workflow. \nolinkurl{BIOMOD_FormatingData} combines presence points, environmental rasters and random pseudo-absences. \nolinkurl{BIOMOD_Modeling} calibrates GLM, GBM, ANN, FDA and MAXNET models under repeated cross-validation. \nolinkurl{BIOMOD_EnsembleModeling} builds mean and coefficient-of-variation ensembles after filtering models with an evaluation threshold. The fitted models are then projected over the current environment and each future climate scenario. A final ensemble forecast combines the corresponding individual projections.

The workflow handles large objects at several stages. Each environmental raster has 63,951,097 cells, and each projection uses five predictor layers. BIOMOD2 may run several model predictions at the same time and may retain or stack their raster outputs. Worker count and projection storage therefore affect the memory schedule even when the scientific parameters do not change.

\section{Code pipeline}

The original pipeline loads the occurrence table and all calibration rasters, selects one species, formats its data, calibrates the individual models and builds the ensembles. It writes model evaluations before projecting the individual and ensemble models over the current environment. A loop then loads each future climate directory, adds the static terrain and soil layers, and repeats both projection stages. The script records one duration for each main phase at the end.

All stages share one R process, one working directory and a set of mutable objects. Species selection, file paths, worker counts and scientific parameters are embedded in the script. There is no scheduler-level isolation between species and no completion marker that distinguishes a finished analysis from a directory left by a failed run.

\section{Original execution times}

\comment{No controlled runtime is available for the unchanged DISTAV script. It cannot run in the repository environment without correcting its Windows paths and syntax, and the preliminary Spartaco records do not identify enough configuration details for a direct comparison.}

\section{Improved execution times}

The measured Leonardo runs use adapted versions of the workflow. Chapter 5 reports those results and keeps them separate from the historical script described here.

\comment{END TODO: Container, workflow and Leonardo}

\section{Preliminary measurements on Spartaco}

Before the Leonardo experiments, Lucia manually recorded four trial runs on Spartaco. The worksheet records per-phase time and RAM observations for two species, with different pseudo-absence settings and CPU counts. Table~\ref{tab:spartaco-preliminary} preserves selected values from those records.

The measurements are useful for reconstructing the starting point of the project. They are not a controlled benchmark: the worksheet does not identify the code revision, container, hardware configuration, input version or measurement method. The source label \texttt{n\_cells} is retained in the table because its meaning cannot be established from the worksheet alone. These results are therefore not compared directly with the later Leonardo measurements.

\begin{table}[htbp]
  \centering
  \caption{Selected fields from four manually recorded trial runs on Spartaco. \texttt{n\_cells} retains the source label. Time units are those recorded in the original worksheet; the maximum RAM is the largest phase value in that row.}
  \label{tab:spartaco-preliminary}
  \resizebox{\textwidth}{!}{%
    \begin{tabular}{llrrrrlll r}
      \hline
      Species & \texttt{n\_cells} & PA & Rep. & RUN & CPU & Formatting & Models & Future projection & Max RAM (GB) \\
      \hline
      \emph{Omalotheca hoppeana} & 681 & 10000 & 10 & 10 & 5 & 1.9 h & 11.33 min & 3.65 min & 52.9 \\
      \emph{Omalotheca hoppeana} & 681 & 10000 & 10 & 10 & 10 & 2.02 h & 7.74 min & 8.07 min & 56.16 \\
      \emph{Agrostis capillaris} & 100000 & 5000 & 5 & 5 & 5 & 32.25 min & 53.87 min & 3.73 min & 26 \\
      \emph{Agrostis capillaris} & 100000 & 10000 & 5 & 5 & 5 & 59.37 min & 1.6 h & 4.22 min & 29.96 \\
      \hline
    \end{tabular}%
  }
\end{table}


% \input{Chapters/RequirementsAnalysis} % Your work
\chapter{My new code}
\comment{Design and implementation of the improvements.}

\section{Container}

\section{Parallel model execution}

\section{Parallel I/O}

% \input{Chapters/Design} % Your work
% \input{Chapters/Implementation} % Your work
\chapter{Experimental results}
\comment{Experimental methodology and results.}

\section{Leonardo}

\section{Tests}

\section{Discussion and lessons learned}

% \input{Chapters/Experiments} % Your work
% \input{Chapters/Discussion} % Your work
\chapter{Conclusions and future development}

\section{Conclusions}

\section{Future development}

\section{Ecoinformatics}
\comment{What it is.}

\comment{Why large datasets are used.}

\comment{Why high-performance computers are used.}

\comment{Purpose.}

\section{Data used}
\comment{In the \texttt{Eval...} matrix, results with \texttt{calibration} below 0.5 are considered less performant.}

\section{Computer science objective}
\comment{To be defined.}

\section{Interdisciplinary collaboration}
\comment{Optional section.}

% \input{Chapters/ConclusionsAndFutureWork}  % Your work
\chapter{Appendix}
\comment{Temporary or optional chapter.}

% \chapter{Appendix: Selected Code Fragments} % Your work
\bibliographystyle{alphaurl}
\bibliography{references}
\end{document}

==== SOME PRECIOUS SUGGESTIONS ON HOW TO WRITE YOUR THESIS ====

The structure heavily depends on the content and type (theoretical, experimental, applicative…).

A Thesis is divided into Chapters. Sometimes, Chapters are grouped into Parts. 

Parts are useful to clearly separate your own contribution from contributions by other scientists. Whether you use Parts or not, it should be always clear to the reader if what you are presenting is needed to provide the context of your work, but has not been developed by yourself. Your own work should be clearly recognizable. The work by others, too.

Chapters are “atomic, consistent and homogeneous units”, dealing with some homogeneous issue. They should be neither too long (in case, split), nor too short (in case, merge). There is no rule stating the right length of a Thesis, and of a Thesis chapter, but - apart from the Introduction and the Conclusions - a 2 pages chapter may seem too short, and a 30 pages one is surely too long. 

A chapter should start with some paragraphs clearly explaining the goal of the chapter, its context, and its main contribution. Sometimes, these few paragraphs introducing the chapter’s contents are presented using a different font to make them more visible.

Sometimes, a Thesis has an Abstract stating, in half a page, the Thesis main goals and achievements. 

Abstract, Introduction and Conclusions are usually written after all the other chapters. 

Do not worry if a chapter is too short/too long, during the writing. Merging and splitting can be faced when the Thesis is almost ready.  

For an internal, “standard” Thesis describing the conception and development of a SW application, a reasonable structure might be (in Latex format):

\input{Chapters/IntroductionAndMotivation} % Your work

\input{Chapters/Background}        % Work done by others!

\input{Chapters/RelatedWork} % Work done by others!

\input{Chapters/RequirementsAnalysis} % Your work

\input{Chapters/Design} % Your work

\input{Chapters/Implementation} % Your work

\input{Chapters/Experiments} % Your work

\input{Chapters/Discussion} % Your work

\input{Chapters/ConclusionsAndFutureWork}  % Your work

\chapter{Appendix: Selected Code Fragments} % Your work

Some chapters can be merged, for example, "Background & Related Work", or "Discussion & Conclusions". Ask your supervisor if this structure may fit your Thesis contents, and how to fill each chapter.

For external Theses, somehow different structures might be adopted, but an external Thesis is, however, a Master Thesis in Computer Science, and it must be accompanied by a good quality document. In particular, for external Theses which are very applicative, we suggest to include a “Lessons Learned” chapter. This chapter, which might substitute or be merged with the “Discussion”, should clarify to the reader which are the main obstacles you met in the practical work you carried out, and the main mistakes that should be avoided by persons facing the same problems you faced. 

If you are working on an external Thesis, ask for further advice to your supervisor, so to avoid to end up with writing a Manual instead of a Thesis. 
Thesis style

This guide is concise and helpful, please read it before starting to write your Thesis: 

https://www.gfdl.noaa.gov/wp-content/uploads/2018/08/Elements_of_Style.pdf

A classical resource for the style of documents is   https://www.chicagomanualofstyle.org/home.html
(trial free access is available).

A good English Dictionary is https://www.merriam-webster.com/

Some facts you cannot ignore, that are based on the experience of your teachers, who collected and listed the most common errors that students use to make. Your supervisor expects that you learn these facts before starting to write the Thesis, and that you carefully check and avoid all the errors explicitly listed in this section. 

Do not underestimate the time and effort required for writing the Thesis. One month before your graduation day, the Thesis should be almost complete: you supervisor and examiner must have the time to check this (almost final) version and give you feedback on it, before uploading it into this AulaWeb module 15 days before your graduation. As a general rule, you might start writing the Thesis three months before your graduation day. If your Thesis is a survey, you should start immediately. If it is a very applicative one, you might start a couple of weeks later. If you know that your English is poor and/or that you spend a lot of time in writing, you should start much earlier.

A Thesis is almost always written in "we" form, even if you are the only author of the document. Some exceptions might be allowed in some chapters, for example when you present the “Lessons Learned” and you comment on your personal experience. Simple rule: use the “we” form and ask your supervisor if you have doubts. 

Never write "We tried to...", "the tool tries to face...", "the Thesis tries to address...", “We hoped to achieve this goal”. Be assertive: "We did", "the tool faces", "the Thesis addresses". The committee will evaluate your work and decide whether you succeeded or not, but by using "try/hope/..." and similar locutions, you convey the very bad feeling that you are not confident in the value of your work. 

On the other hand, never write enthusiastic sentences like "This tool has great performances...", "The results we got are great", “The code is so easy to read”, etc. Your statements must be objective and backed up by evidence. You can write that your tool has some performances and that these performances are better than those of toolX, toolY, toolZ. Of course, you must prove what you write. You can write that you developed 1000 Lines of Code (LOC), but you cannot state that it is “easy to read” unless you are carrying out experiments on code readability, and you can measure this feature in a scientific way.

The committee does not care if you are satisfied: the scientific content of the Thesis is the only thing that matters, and that will be evaluated. So, avoid sentences like "We are very happy with the tool we developed..."

Use short and simple sentences.

Use as few adverbs and adverbial locutions as possible. Most of them are completely useless, and just make the reading harder. Try to remove them and check the result. If the meaning of the sentence is not affected, than definitively remove them.

Do not use slang and abbreviations: the Thesis is a scientific document and must meet a scientific style. Do not use "Don't", "Can't", "Let's". Use "Do not", "Can not/cannot", "Let us".

All figures and tables must have meaningful, informative captions. Tagging all of them as “Experiment” conveys no information. Also, they must be referenced from inside the text and explained in a satisfactory way in your Thesis. 

Plots on the Cartesian plane must include textual, informative labels for x- and y- axes; in general, all kinds of graphs (bar charts, histograms, pies) must include labels to allow the reader to immediately grasp their content.  

If you use or define an equation, all the terms which appear inside it must be properly introduced and explained in the text, before the equation itself. 

All captions and footnotes must start with uppercase. 

The words Section/Figure/Table/ etc must be uppercase when followed by a number ("In Section 2.4 we present..."), lowercase in the other cases ("In this section, we address...").

All the figures, tables, bibliographic items must be referenced from inside the text.

Bibliography: all the entries must be as complete as possible, even those for web sites (you cannot just put a url). You find already made, complete bibliographic items for articles and books in many existing databases such as DBLP (https://dblp.uni-trier.de/) and ACM.

If you use Latex, you may safely use the BibTex items in both DBLP and the ACM library. Do not fully trust the BibTex items you find in CiteSeer and Google Scholar: sometimes their quality is very low. 

The months you spend on your Thesis are special ones: you can do what you always dreamt of, you can experiment, you can be autonomous in your choices, but still counting on very experienced professors willing to help you.

And hence … stick to the guidelines and enjoy your Thesis!